\section{Machine learning: what was tried, and what it gave}
\label{sec:ml}

This section is an accounting, not an advertisement. None of the theorems in
this paper were discovered by a neural model. They came from exhaustive
enumeration, LP duality, and machine-checked mining over finite catalogs, as
described in the preceding sections. We nonetheless ran a substantial
interpretability programme alongside that work, on the hypothesis that a
small model trained to predict D-reducibility might expose a human-statable
sufficient condition that the enumeration was not organized to find. We
report what that programme gave: one modest positive result, one
methodological finding about the fragility of ``mechanism'' as a concept,
and several negative results we think are worth recording so that others do
not have to rediscover them. We write in the register of a referee report
on our own work.

\subsection{Design and motivation}

The programme deliberately transplanted the template of
Huang, Jackson, and Lee's mechanistic-interpretability study of the zeta
map \cite{huang2025zeta}: train a small transformer on a well-defined
mathematical computation, probe its internal representations, and try to
extract a human-readable mechanism. That template sits inside the broader
agenda of using machine learning to guide mathematical intuition
\cite{davies2021}. The target computation here was D-reducibility itself:
given a configuration, predict the Kempe-closure trace produced by the
D-reduction search and the final verdict (reducible or not). The rationale
was simple. If a small model can predict reducibility reliably, whatever
internal feature it relies on is a candidate for a sufficient condition a
human could then state and verify independently -- the same logic that
motivated the zeta-map study's search for a legible mechanism.

We built two model generations.

\textbf{D1.} An encoder-decoder transformer ($d_{\mathrm{model}}=192$,
4 heads, feed-forward width 384, 2 encoder layers, 2 decoder layers,
dropout 0, 1{,}543{,}725 parameters), trained seq2seq from a canonical
token sequence encoding a configuration to a target sequence of
closure-trace count buckets followed by a verdict token. Training ran 120
epochs at learning rate $10^{-3}$, batch size 64, seed 0, for 10{,}392~s on
an Apple MPS device. The corpus, \art{data/d1\_corpus.jsonl}, held
14{,}036 deduplicated records spanning ring sizes 6--16 (633 from RSST,
2{,}614 from Steinberger, 5{,}961 generated, 4{,}828 drawn from the 2026
pool); after excluding records with degenerate traces or out-of-range
length, 6{,}559 remained, split 5{,}904 train / 655 validation. Artifacts:
\art{checkpoints/d1\_full\_v1.pt},
\art{results/d1/train\_metrics\_full\_v1.json},
\art{results/d1/corpus\_stats.json}.

\textbf{D1-v2.} A second generation of per-ring specialists at
$r = 8, 9, 10$, sharing a 2-layer encoder ($d_{\mathrm{model}}=256$, 4
heads, feed-forward width 512, masked-mean-pool summary) with two
alternative heads. A ``dynamics'' head iterates a weight-tied
\texttt{GRUCell} over the per-ring survivor-set bitmap with per-round
supervision (3.26M parameters at $r=9$, 6.64M at $r=10$); a
``verdict-only'' \textsc{control} head reads the verdict directly off the
encoder summary with no recurrence (1.19M parameters), and exists
specifically to test whether the dynamics head's recurrent machinery
contributes anything the summary vector does not already contain. Data
lived in \art{data/v2/traces\_r8.jsonl}, \art{data/v2/traces\_r9.jsonl},
and \art{data/v2/traces\_r10.jsonl}, carrying full per-round survivor sets;
the corpora held 1{,}162 / 1{,}446 / 1{,}746 configurations at $r=8,9,10$, with
ring-colouring code-index sizes 1{,}094 / 3{,}281 / 9{,}842 respectively. Artifacts: \art{results/d1v2/r9/},
\art{results/d1v2/r10/}, \art{results/d1v2/r10v2/}.

\subsection{Results}

D1 reaches verdict accuracy 0.9298 on its 655-example validation split,
against a majority-class baseline of 0.5603 -- a real signal, well above
chance. Full-sequence exact match (verdict token plus the entire
closure-trace bucket sequence) is a much weaker 0.2397: the model gets the
answer right far more often than it gets the reasoning trace right.

We then asked what the encoder actually represents, using a probe ladder
recorded in \art{results/d1/interp/report.md} (ridge logistic probes on
standardized inputs, probe-test $n=131$, 95\% bootstrap confidence
intervals with \texttt{n\_boot}$=2000$). A 16-feature hand-coded shallow
baseline -- counts and simple graph statistics of the configuration, no
learned representation involved -- already reaches 0.8779
$[0.8168, 0.9313]$. The best pure-encoder representation is
\texttt{encoder\_layer0\_mean}, the \emph{pre-attention} embedding bag, at
0.9008 $[0.8473, 0.9466]$: a delta of $+0.0229$ over the shallow baseline,
but with overlapping confidence intervals, so not a distinguishable
improvement. The post-attention layers are worse, not better: layers 1 and
2 both give 0.8855, despite fitting the training set considerably better
than layer 0 (train accuracy 0.9599 and 0.9866 respectively, against 0.8817
for the pre-attention bag). Attention is improving the model's memorization
of the training distribution without improving the generalizing signal a
downstream probe can extract. Augmenting the shallow features with
hand-derived structural candidates reaches 0.9237
$[0.8779, 0.9695]$ -- that is, hand features beat the encoder's best
representation outright. We state this plainly as the honest summary of
D1: the transformer's internal representation of reducibility is, on this
evidence, no better than what a referee could write down with a pencil,
and in places worse.

A second question was whether the decoder accumulates verdict evidence as
it emits the closure trace, the way one might hope a legible mechanism
would. We probed decoder hidden states at each greedy decoding step
$t = 0, \dots, 39$. Verdict-probe accuracy is 0.8473 at $t=0$, before a
single trace token has been produced, and the curve stays essentially flat
thereafter: the mean over $t=0$--$2$ is 0.8702, the mean over $t=35$--$39$
is 0.8656, and every later confidence interval overlaps the one at $t=0$.
The closure-trace tokens ride along with a verdict decision already
implicit at the start; they do not build toward one.

\subsection{The boundary result}
\label{sec:ml-boundary}

The results above concern accuracy on the full validation distribution,
which is dominated by easy cases. We pre-registered a harder ``boundary''
stratum, flagged by \texttt{tools/d1v2\_datagen.py::boundary\_flag}, to test
whether any representation carries information beyond the shallow features
specifically on the cases an easy classifier would get wrong: a
configuration is boundary if it is non-reducible with
$0 < n_{\mathrm{consistent}} \le 50$, or reducible with rounds $\ge 6$. We
pre-registered the decision rule before looking at boundary results: a
representation beats the shallow baseline only if its boundary 95\%
bootstrap confidence interval does not overlap the shallow baseline's and
lies strictly above it.

The first pass, at rings 9 and 10 on the smaller corpus, was a clean null:
0 of the 7 non-baseline representations at $r=9$ and 0 of 7 at $r=10$ cleared the
rule, and
the boundary strata in the probe-test splits were tiny (7 and 11 examples
respectively) -- underpowered by construction.

The second pass used an enlarged ring-10 corpus (3{,}269 configurations;
probe-test $n=98$, of which 21 are boundary) and gave a real result.
Table~\ref{tab:boundary} summarizes it.

\begin{table}[htbp]
\centering
\begin{tabular}{lcc}
\toprule
Representation & Overall accuracy & Boundary accuracy [95\% CI] \\
\midrule
shallow & 0.7551 & 0.5238 $[0.3321, 0.7143]$ \\
shallow\_plus\_structural & 0.8571 & 0.7619 $[0.5714, 0.9048]$ \\
control\_encoder\_summary & 0.9082 & 0.9048 $[0.7619, 1.0000]$ \\
dynamics\_encoder\_summary & -- & 0.9048 \\
dynamics\_h0 & -- & 0.9048 \\
dynamics\_h2 & -- & 0.9048 \\
\bottomrule
\end{tabular}
\caption{Ring-10 boundary probe results, second pass ($n=98$ overall,
21 boundary). \texttt{control\_encoder\_summary} and three dynamics-model
representations clear the pre-registered rule against the shallow
baseline.}
\label{tab:boundary}
\end{table}

\texttt{control\_encoder\_summary} -- the 256-dimensional encoder summary
feeding the verdict-only control head, with no recurrent dynamics involved
at all -- clears the pre-registered rule against the plain shallow
baseline, as do three dynamics-model representations. This is the one
genuinely positive result in the programme: some encoder representation
carries boundary-relevant information the 16 hand-picked shallow features
do not.

We caveat it on three counts. First, $n=21$ boundary examples is a small
stratum; the confidence intervals in Table~\ref{tab:boundary} are wide.
Second, against the also-hand-engineered \texttt{shallow\_plus\_structural}
baseline the comparison is never confidence-interval-decisive at any seed
we tried (0.9048 versus 0.7619, intervals overlap); the result holds
against the naive 16-feature baseline, not against a reasonably strengthened
one. Third, we note for the record that the source report's own concluding
prose (\art{results/d1v2/v2run/interrogation.md}) contradicts its table: it
recycles wording from the earlier null run at rings 9 and 10. We report the
table, not the prose.

\subsection{Extraction: hand features recover most, but not all, of the advantage}
\label{sec:ml-extraction}

Given that some encoder representation beats the shallow baseline on
boundary cases, we tried to extract what it was finding, in the spirit of
turning a learned feature into a stated one. We hand-coded 14 arrangement
features of $H_5$, the subgraph induced on interior vertices of degree
exactly 5 (\art{src/fourcolor/d1\_features.py}): number of connected
components, largest component size, number of isolated vertices, edge
count, density, maximum degree, existence of a length-3 path, number of
vertices adjacent to the ring, minimum and mean component-to-ring
distance, minimum and mean inter-component distance, the count of
degree-5 vertices with at least 3 degree-5 neighbours, and the existence
of a diamond in $H_5$.
Appending these 14 to the 24-dimensional \texttt{shallow\_plus\_structural} vector
(16 shallow features plus 8 structural ones) gives a 38-dimensional
\texttt{shallow\_plus\_arrangement} vector. Table~\ref{tab:extraction}
reports boundary accuracy across four seeds, drawn from
\art{results/d1v2/extraction/report.md}.

\begin{table}[htbp]
\centering
\begin{tabular}{lccccc}
\toprule
Seed & $n_{\mathrm{boundary}}$ & shallow & +structural & +arrangement [95\% CI] & control\_encoder \\
\midrule
0 & 21 & 0.5238 & 0.7619 & 0.9524 $[0.8571, 1.0000]$ & 0.9048 \\
1 & 23 & 0.4348 & 0.7826 & 0.9130 $[0.7826, 1.0000]$ & 0.7826 \\
2 & 24 & 0.5417 & 0.7083 & 0.8750 $[0.7083, 1.0000]$ & 0.7500 \\
3 & 25 & 0.5200 & 0.7600 & 0.7600 $[0.5990, 0.9200]$ & 0.9200 \\
\bottomrule
\end{tabular}
\caption{Boundary accuracy of the hand-coded arrangement-feature extraction
against the plain shallow baseline, the structural baseline, and the
encoder control, across four seeds.}
\label{tab:extraction}
\end{table}

The extraction succeeds on three of four folds: at seeds 0, 1, and 2 the
arrangement features push boundary accuracy well above
\texttt{shallow\_plus\_structural}, and in all three above the encoder control as
well (0.9524 vs 0.9048, 0.9130 vs 0.7826, 0.8750 vs 0.7500). It fails on seed 3: the arrangement features add
exactly nothing over the structural baseline there, an identical point
estimate of 0.7600 to 0.7600, while the encoder control reaches 0.9200. We
do not know why seed 3 is different; we simply report that it is.

The boundary result of Table~\ref{tab:boundary} itself replicates across
these same four seeds in direction but not always in confidence-interval
strength: the control encoder beats the plain shallow baseline by $+0.38$,
$+0.35$, $+0.21$, and $+0.40$ at seeds 0--3 respectively, so the sign is
consistent 4 of 4. But the confidence-interval-decisive version of the
comparison -- the pre-registered rule -- holds only at seeds 0 and 3, i.e.\
2 of 4. The headline boundary result reported in
of \S\ref{sec:ml-boundary} was therefore somewhat fortunate in its seed choice;
a differently seeded first pass could plausibly have reported an
inconclusive result instead. Artifacts:
\art{results/d1v2/extraction/report.md},
\art{results/d1v2/replication/report.md}.

We also performed a small error-set digestion at seed 0
(\art{results/d1v2/digestion/cases.md}). Of the 21 boundary probe-test
examples, the encoder control is right and the structural baseline is
wrong on 4; the reverse holds on 1; both are right on 15; both are wrong
on 1. All 4 configurations where the encoder wins have $n=18$, are
D-reducible, and have closure round counts $6, 6, 7, 6$. Their mean count
of degree-5 diamonds is 2.75, against 0.88 for the rest of the validation
boundary set and 0.97 corpus-wide -- roughly a factor of three -- with
every other feature we measured flat across the two groups. We decline to
run a significance test at $n=4$; the sample is too small to support one
honestly.

We flag one honest gap here. The automated per-configuration check that
the arrangement features specifically fix these four cases did not
actually run: \art{tools/d1v2\_extraction\_test.py} guards its
\texttt{error\_set\_check} behind
\texttt{hasattr(interro, "probe\_predictions")}, which is \texttt{False}
for the interrogation object actually produced, so the check is silently
skipped and \texttt{error\_set\_check} is empty at every seed. The causal
link we would like to draw -- diamond count explains the four flipped
cases -- is therefore suggestive, resting on the aggregate statistic above,
and not established at the level of individual predictions.

\subsection{Mechanism is training-contingent: a replication finding}

This subsection reports a companion experiment run in a separate
repository, \texttt{/Users/yugendren/experiments/zeta\_map\_interp}, not in
the codebase for this paper. Its subject is the published zeta-map result
we took as our template, \cite{huang2025zeta}, and the question is whether
its reported mechanism replicates under a faithful reimplementation.

We reimplemented the published architecture (1 encoder layer, 1 decoder
layer, 1 attention head, $d_{\mathrm{model}}=128$, feed-forward width 256,
post-norm, learned positional embeddings, weight-tied output projection):
337{,}029 parameters against the paper's reported 339{,}716. Training
reached the same headline accuracy the paper reports: 0.9950 greedy
exact-match on a held-out instances of size $n=13$. We then reran each of the
paper's three reported mechanistic findings against our own checkpoint.

\begin{itemize}
\item \emph{Causal ablation.} The paper reports that ablating
cross-attention to up-steps costs about 1.5\% exact-match accuracy, read
as evidence of a selective mechanism attending preferentially to up-steps.
We measure a baseline of 0.996 exact match falling to 0.0000 when
ablating attention to up-steps, and to 0.0005 when ablating attention to
the complementary down-steps. Both ablations are catastrophic and
symmetric: cross-attention is load-bearing wholesale in our model, with no
selectivity toward either step type.

\item \emph{Attention statistics.} Mean cross-attention mass on up-steps
is 0.5344 against a uniform baseline of 0.5; the attention argmax lands on
an up-step for 0.5739 of queries; only 0.0682 of queries have their argmax
at the path's own maximum level; the mean level of the argmax-attended
position is 2.369, against 2.471 averaged over all source positions. These
are essentially uniform statistics -- we find no evidence of attention
selecting the path's highest level.

\item \emph{Linear probes.} The paper reports approximately 84\% probe
accuracy for recovering path levels from encoder hidden states. We measure
0.9478 (14 classes, majority baseline 0.2439) for level, and 1.0000 for
step type (2 classes, baseline 0.5). This one direction replicates, and in
fact strengthens.

\item \emph{Length generalization.} The paper evaluates lengths
$n = 11$--$16$. Our model scores exactly 0.0000 exact match at
$n = 11, 12, 14, 15, 16$, with per-token character accuracy 0.5872 at
$n=12$ against roughly 0.5 chance. We tried four positional-placement
variants at evaluation time (left-aligned to length 24, left-aligned
padded to length 26, right-aligned, centered); all four give 0.0000 exact
match, which rules out a pure positional-embedding-range explanation for
the failure.
\end{itemize}

We draw the conclusion carefully, and without asserting that the original
authors made an error: same architecture, same task, the same held-out
accuracy we would need to call the reimplementation faithful, and yet a
materially different internal mechanism by every measure we could apply.
Whatever mechanism a small transformer settles on for a mathematical
computation appears to be, at least in part, a property of the particular
training run rather than a property of the task it was trained on. For an
interpretability-driven discovery programme this is the central
methodological risk: an extracted ``mechanism'' may describe a training
run rather than a fact about the mathematics. This is exactly the failure
mode we then observed inside our own setting, at seed 3 of
\S\ref{sec:ml-extraction}.

\subsection{Assessment}

None of the four contributions of this paper came from the neural models.
The one thing the machine-learning line delivered that fed forward into
the rest of the paper was the degree-5 arrangement feature family of
\S\ref{sec:ml-extraction}, which sharpened our attention on $H_5$ and became
one input, among others, to the rule-mining process that produced the
candidate rules discussed in Section~\ref{sec:negative} -- of which the
surviving rule, in its exhaustively verified form, became
Theorem~\ref{thm:fr}. We describe that as an indirect, heuristic
contribution and make no stronger claim than that.

We also record two structural failures of the programme, so that the
positive results above are read in proportion. First, the dynamics
model -- the entire interpretability design bet, and the closest analogue
to the zeta-map paper's step-by-step mechanism -- attains teacher-forced
verdict accuracy 1.0000, but collapses under free-running rollout to
0.7778 at $r=9$, 0.7786 at $r=10$, and 0.6245 at $r=10$ in the second
($r10v2$) configuration, with free-running boundary accuracy exactly
0.0000 at $r=10$ and $r10v2$: unrolled without teacher forcing, the model
converges to the trivial fixed point of predicting that nothing is ever
eliminated. Second, the D1 checkpoint's learned positional embeddings
extend only to length 259, which admits 7{,}021 of the 14{,}036
configurations in the corpus; ring sizes above roughly 13 were structurally
out of the model's reach, independent of anything about the underlying
mathematics.

Finally, we caught and quarantined one tautological feature during this
work. The quantity $n_{\mathrm{consistent}}$ correlates with the verdict
label at Spearman $-0.850$ and appears to ``close 167\% of the gap'' to
perfect boundary accuracy -- because, on inspection, it is the label in
disguise: $n_{\mathrm{consistent}} \le 50$ is one of the two disjuncts
defining boundary non-reducibility in the first place. It is excluded from
consideration via \texttt{TAUTOLOGICAL\_CANDIDATES} in
\art{tools/d1\_interrogate.py}, and we mention it here only as a caution to anyone repeating this
kind of search: a feature that predicts the label unusually well is worth
checking for circularity before it is reported as a discovery.
